% 1.4 点群
\section{点群}

三维点群是作用在点 $O$ 上（因而保持 $O$ 不动）、并使三维欧氏空间 $E(3)$ 中所有距离与角度都保持不变的对称算符群。具有这些性质的对称算符是作用在 $O$ 点的幺正算符，即绕过 $O$ 点的轴的转动，或这类转动与反演的乘积；这类乘积包括过 $O$ 点的平面内的反射。

若该群包含所有可能的转动而没有其他元素，则称为\textit{三维真转动群}（3-dimensional proper rotation group），它与由行列式为 $+1$ 的全部 $(3 \times 3)$ 正交矩阵构成的群 $SO(3)$ 同构。理由是：在转动 $R$ 作用下，位置矢量为 $\mathbf{r}$ 的点 $P$ 被移到位置矢量为 $\mathbf{r}'$ 的点 $P'$，

\begin{equation*}
\mathbf{r}' = \mathbf{R}\mathbf{r},
\tag{1.4.1}
\end{equation*}

其中 $\mathbf{R}$ 是 $(3 \times 3)$ 矩阵。由于所有距离与角度都必须保持不变，故对一切可能的向量对 $\mathbf{r}$、$\mathbf{s}$ 都有

\begin{equation*}
\mathbf{r}.\mathbf{s} = \mathbf{Rr}.\mathbf{Rs} = \mathbf{r}'.\mathbf{s}',
\tag{1.4.2}
\end{equation*}

而这成立当且仅当 $\mathbf{R}$ 是正交矩阵。众所周知，正交矩阵的行列式为 $+1$ 或 $-1$；实际上只有行列式为 $+1$ 的正交矩阵才对应转动，因为零角度转动是恒等操作、其行列式为 $+1$，再由连续性要求，任何对应转动的矩阵行列式也都是 $+1$。

若该群包含所有可能的转动及其与反演的乘积，则称为\textit{三维转动群}（3-dimensional rotation group），它与全部 $(3 \times 3)$ 正交矩阵构成的群 $O(3)$ 同构。对应于行列式为 $-1$ 的矩阵的算符称为非真转动（improper rotations），它们是真转动与反演的乘积。反演与所有转动对易：这由它在 $O(3)$ 中表示为负单位矩阵立即可得。

这里必须明确我们处理幺正算符的约定。约定是：我们把幺正算符理解为\textit{作用在空间的点上}，而这些点的坐标总是相对一组固定坐标轴给出。移动空间的点而保持坐标轴不动的算符称为\textit{主动}（active）算符；移动坐标轴而保持点不动的算符称为\textit{被动}（passive）算符。\textit{本书一律采用主动约定。}联系两种约定的公式是

\begin{equation*}
\mathbf{R}_{\mathrm{active}} = \mathbf{R}_{\mathrm{passive}}^{-1}.
\tag{1.4.3}
\end{equation*}

从一种约定换到另一种时必须非常小心，因为两种约定的对应不是自同构而是反自同构：也就是说，若 $\mathbf{R}$、$\mathbf{S}$、$\mathbf{T}$ 是三个转动的矩阵，且在主动模式下有

\begin{equation*}
\mathbf{R}_{\mathrm{active}}\mathbf{S}_{\mathrm{active}} = \mathbf{T}_{\mathrm{active}}
\tag{1.4.4}
\end{equation*}

则被动算符的相应公式是

\begin{equation*}
\mathbf{S}_{\mathrm{passive}}\mathbf{R}_{\mathrm{passive}} = \mathbf{T}_{\mathrm{passive}},
\tag{1.4.5}
\end{equation*}

即左边算符的乘法次序要颠倒，而不是像自同构那样保持不变。

三维转动群的每个子群都是一个\textit{点群}（point group）。真转动群的子群称为\textit{真点群}（proper point groups）。真点群只包含绕过 $O$ 的轴的转动。以下只讨论有限阶点群。有限阶真点群可分类如下：

(i) \textit{循环群}. 当只有一条转动轴（比如 $Oz$）时出现循环群。国际记号记作 $n$，Schönflies 记号记作 $C_{n}$。对每个正整数 $n$ 都有一个明确的循环群，其 $n$ 个元素是绕该轴转 $2\pi q/n$ 角的转动，$q = 0, 1, 2, \ldots, (n-1)$。在 Schönflies 记号中，绕 $Oz$ 轴的 2$\pi/n$ 逆时针转动记作 $C_{nz}^{+}$，相应的顺时针转动记作 $C_{nz}^{-}$；当 $n = 2$ 时上标 $\pm$ 省略。

(ii) \textit{二面体群}. 由把正 $n$ 棱柱变为其自身的全部转动组成。国际记号：$n$ 为偶数时记 $n22$，$n$ 为奇数时记 $n2$；Schönflies 记号 $D_{n}$。对每个正整数 $n \geqslant 2$ 有一个明确的二面体群，其 $2n$ 个元素是：绕棱柱棱方向主轴转 $2\pi q/n$ 的 $n$ 个转动（该轴称为\textit{主轴}），以及 $n$ 个阶为 2 的转动（即转过 $\pi$ 角）——每个这样的转动绕一条\textit{副轴}（secondary axis），副轴垂直于主轴。共有 $n$ 条副轴，对称地布置在棱柱内，相邻两条副轴的夹角为 $\pi/n$。$n$ 为偶数时一半副轴穿过棱柱面的中心，另一半穿过棱的中点；$n$ 为奇数时每条副轴既穿过一个面的中心又穿过对边的中点。在 Schönflies 记号中，绕主轴的 $n$ 个转动的标号与循环群相同，绕副轴 $Op$ 的转动记作 $C_{2p}$。

(iii) \textit{四面体群}. 由把正四面体变为其自身的全部转动组成。国际记号 23，Schönflies 记号 $T$。该群有四条过四面体顶点的 3 阶轴，以及三条连接对边中点的 2 阶轴。三条 2 阶轴构成一个直角三联组，因此方便地记作 $Ox$、$Oy$、$Oz$。若把四面体的顶点记作 1、2、3、4，则四条 3 阶轴记作 01、02、03、04（见图 1.3）。按上述 Schönflies 记号，23（$T$）的 12 个元素为：$E$（单位元），$C_{3j}^{\pm}$（$j = 1, 2, 3, 4$），以及 $C_{2m}$（$m = x, y, z$）。

(iv) \textit{八面体群}. 由把立方体变为其自身的全部纯转动组成。国际记号 432（有时简写为 43），Schönflies 记号 $O$。正四面体 1234 可以内接于立方体，使其四个顶点位于立方体的四个角上；这样便可把立方群的元素与四面体群的元素联系起来——事实上 23（$T$）是 432（$O$）的子群。轴 01、02、03、04 仍是 3 阶轴，但此时 $Ox$、$Oy$、$Oz$ 是 4 阶轴。此外还有六条连接立方体对边中点的 2 阶轴，按图 1.3 记作 $Oa$、$Ob$、$Oc$、$Od$、$Oe$、$Of$。于是 432（$O$）的 24 个元素为：$E$，$C_{3j}^{\pm}$（$j = 1, 2, 3, 4$），$C_{4m}^{\pm}$（$m = x, y, z$），$C_{2m}$（$m = x, y, z$），以及 $C_{2p}$（$p = a, b, c, d, e, f$）。

(v) \textit{二十面体群}. 由把正五角十二面体或正二十面体变为其自身的全部转动组成。它有六条 5 阶轴、十条 3 阶轴、十五条 2 阶轴，因此共有 60 个元素（Cohan 1958，Cotton 1963）。国际记号 532，Schönflies 记号 $Y$。

若用 $\bar{1}$（或 $C_{i}$）表示由单位元 $E$ 和反演 $I$ 组成的群，则对任一真点群 $\mathbf{P}$，总可以按定义 1.2.13 构造外直积 $\mathbf{P} \otimes \bar{1}$（或 $\mathbf{P} \otimes C_{i}$）。这样就得到一族新的、为 $O(3)$ 子群的点群。这些点群（我们给出国际记号，括号内为其 Schönflies 等价记号）如下：

\begin{quote}
(i) 由 $n$（$C_{n}$）（$n$ 为奇数）及其不变子群可得 $\bar{n}$（$S_{2n}$）；由 $n$（$C_{n}$）（$n$ 为偶数）可得 $n/m$（$C_{nh}$）。（在 Schönflies 记号中 $S_{2}$ 常被 $C_{i}$ 代替。）
(ii) 由 $n2$（$D_{n}$）（$n$ 为奇数）可得 $\bar{n}m$（$D_{nd}$）；由 $n22$（$D_{n}$）（$n$ 为偶数）可得 $n/mmm$（$D_{nh}$）。（当 $n = 2$ 时国际符号通常写作 $mmm$。）
(iii) 由 23（$T$）可得 $m3$（$T_{h}$）。
(iv) 由 432（$O$）可得 $m3m$（$O_{h}$）。
(v) 由 532（$Y$）可得 $\bar{5}3m$（$Y_{h}$）。
\end{quote}

由真点群构造新点群还有第二种方法：当 $\mathbf{P}$ 有指数为 2 的不变子群 $\mathbf{Q}$ 时特别有效。此时可写

\begin{equation*}
\mathbf{P} = \mathbf{Q} + R\mathbf{Q}
\tag{1.4.6}
\end{equation*}

其中 $R$ 是某个转动。于是群

\begin{equation*}
\mathbf{P}' = \mathbf{Q} + IR\mathbf{Q},
\tag{1.4.7}
\end{equation*}

其中 $I$ 是反演，也是一个点群。用这种方法导出的点群如下：

\begin{quote}
(i) 由 $2n$（$C_{2n}$）（$n$ 为奇数）及其不变子群 $n$（$C_{n}$）可得 $\bar{2}n$（$C_{nh}$）。（其中 $\bar{2}$（$C_{1h}$）通常被记号 $m$（$C_{s}$）代替。）由 $2n$（$C_{2n}$）（$n$ 为偶数）及其不变子群 $n$（$C_{n}$）可得 $\bar{2}n$（$S_{2n}$）。
(ii) 由 $n22$ 或 $n2$（$D_{n}$）及其不变子群 $n$（$C_{n}$）可得 $nmm$ 或 $nm$（$C_{nv}$）；由 $(2n)2$（$D_{2n}$）（$n \geqslant 2$）及其不变子群 $n22$ 或 $n2$（$D_{n}$）可得 $(\bar{2}n)2m$（$D_{nh}$）（$n$ 为奇数）与 $(\bar{2}n)2m$（$D_{nd}$）（$n$ 为偶数）。
(iii) 23（$T$）没有指数为 2 的不变子群，因此不能用这种方法从 23（$T$）导出新群。
(iv) 由 432（$O$）及其不变子群 23（$T$）可得 $\bar{4}3m$（$T_{d}$）。
(v) 532（$Y$）没有不变子群，因此不能用这种方法从 532（$Y$）导出新群。
\end{quote}

这两种由真点群构造新点群的方法穷尽了 $O(3)$ 的一切有限阶的、不纯由真转动组成的点群。有限点群的分类至此完成。读者应该已经注意到，点群及其对称元素的标记至少有两套并行的体系：国际记号与 Schönflies 记号。本书对点群与空间群同时使用两种记号：Schönflies 记号总是放在国际记号后的括号里。

现在稍稍展开一下这些记号的含义。国际记号中，数字 $n$ 表示存在一条 $n$ 阶转动轴；数字上的横杠 $\bar{n}$ 表示存在一条 $n$ 阶旋转反演轴；字母 $m$ 表示存在一个反射面。此外，若字母 $m$ 与数字之间用斜线（/）隔开，则表示有一个对称面垂直于主转动轴。最后，点群的 Schönflies 记号中下标 $h$、$v$、$d$ 分别表示水平、竖直与对角的反射面——这只有在点群处于某种标准定向时才有意义。很多情况下我们需要使用非标准定向的点群，这也是我们更愿意用国际记号来标记点群本身的原因。不过这不会引起混淆，因为在非标准定向下每个元素都会得到一个确定的标记，读者可对照表 1.3 辨认。

本节开头我们说过，点群是作用在一点 $O$ 上、并使三维欧氏空间 $E(3)$ 中所有距离与角度保持不变的对称操作群。本书关心的是晶体对称操作群的表示理论，因此讨论将限于\textit{晶体学点群}（crystallographic point groups）。晶体学点群必须满足附加条件：与某种晶体的平移对称性相容。可能满足该条件的点群操作是：恒等操作、反演、某些平面内的反射，以及绕某些 1、2、3、4 或 6 阶轴的转动。可以证明，由这些操作本质上不同的组合只有有限多种；事实上晶体学点群恰有 32 个，许多书都讨论了它们的各种性质（例如 Buerger (1956)，Koster, Dimmock, Wheeler, and Statz (1963)，Phillips (1963\textsuperscript{a})）。32 个点群可以按主轴的阶归入七个\textit{晶系}（crystal systems，有时称 syngonies）：主轴为 1、2、3、4 或 6 阶的单一主轴点群分属三斜、单斜、三方、四方、六方五个晶系；另有正交晶系（有三条互相垂直的 2 阶轴）与立方晶系（有四条指向正四面体顶点的 3 阶轴），合为七系。本章描述的点群大多数是晶体学点群；非晶体学点群的例子有二十面体群以及阶数 $n$ 不取 1、2、3、4、6 的循环群和二面体群。

\textit{点群操作的记号}

文献中点群对称元素的标记方式非常多，没有一种是完全令人满意的（各种记号的讨论见 Warren (1968)）。有些作者（例如 Kovalev (1965)，Miller and Love (1967)，Slater (1965\textsuperscript{b})）苦于记号不统一，干脆给立方群 $m3m$（$O_{h}$）的元素编号 $\mathbf{R}_{1}, \mathbf{R}_{2}, \ldots, \mathbf{R}_{48}$，给六方群 $6/mmm$（$D_{6h}$）编号 $\mathbf{R}_{1}, \mathbf{R}_{2}, \ldots, \mathbf{R}_{24}$；这种方案在使用计算机时非常方便。然而，成熟的记号虽然不完美，却带有一定的含义（任意编号则没有），因此本书对点群与空间群的实际对称操作采用经我们扩充的 Schönflies 记号，使每个点群算符都能被单独辨认。真转动的记号已在真点群的分类中给出；完整的辨认由表 1.2、表 1.3 以及图 1.1--1.3 给出。

\begin{figure}[H]
\centering
\includegraphics[width=0.60\textwidth]{fig1_1.png}
\caption{三斜、单斜、正交和四方晶系的对称元素。这些晶系中的点群都是 $m3m$（$O_{h}$）的子群，故采用相同记号。$x$、$y$、$z$ 构成右手系。各对称操作的记号标在图中该操作把字母 $E$ 移到的位置上。}
\end{figure}

\begin{figure}[H]
\centering
\includegraphics[width=0.62\textwidth]{fig1_2.png}
\caption{三方和六方晶系的对称元素。$z$ 轴垂直于纸面向上。各对称操作的记号标在图中该操作把字母 $E$ 移到的位置上。}
\end{figure}

许多结果的有用性依赖于对称元素的完整辨认，因此我们对此略加讨论。对属于大多数晶系的点群，可以用极射赤面投影图（stereogram）之类的平面图来辨认对称元素：图 1.1 用于三斜、单斜、正交和四方晶系，图 1.2 用于三方和六方晶系。Schiff（1954）曾用双面纸模型帮助想象各种非立方点群。而对立方晶系，直接从立方体图形辨认对称元素更容易，见图 1.3。$C_{nr}^{+}$ 表示把空间的点绕图中标记为 $r$ 的轴逆时针转 $2\pi/n$ 弧度的操作，$C_{nr}^{-}$ 表示绕同一轴顺时针转 $2\pi/n$ 弧度的操作；$E$ 是恒等操作，$I$ 是反演。表 1.2 按七个晶系统一列出了各晶系中出现的算符；表的每一部分分两列，右列元素是左列相应元素与 $I$ 的乘积。表中给出的是相应点群处于相对图的 $x$、$y$、$z$ 轴的标准定向时的元素。某些点群有时必须使用非标准定向，但由于每个元素都有确定的标记，这不会引起混淆。应当注意：反射面一律记作 $\sigma$，其他转动反射记作 $S$；还应指出，在我们采用的记号里 $IC_{2} = \sigma$、$IC_{3}^{+} = S_{6}^{-}$、$IC_{4}^{+} = S_{4}^{-}$、$IC_{6}^{+} = S_{3}^{-}$（即 $IC_{3}^{+} = S_{6}^{-}$，等等）。这是因为 $S_{n}^{+}$ 表示绕轴逆时针转 $2\pi/n$ 后再在垂直于该轴的平面内反射，于是 $IC_{n}^{+} = \sigma C_{2}C_{n}^{+} = \sigma(C_{n}^{+})^{n+2} = (S_{2n}^{-})^{n+2}$。

还有一个重要的、前面已经提过的要点必须强调：处理点群算符时，我们始终把它理解为\textit{主动}算符——移动空间的点而保持坐标轴不动；同样，写成 $\{R \,|\, \mathbf{v}\}$ 形式的空间群算符（见 1.5 节）也按主动算符解释。

\begin{figure}[H]
\centering
\includegraphics[width=0.58\textwidth]{fig1_3_scan.png}
\caption{立方晶系的对称元素（Altmann and Bradley 1963\textsuperscript{b}）。}
\end{figure}

表 1.3 给出 32 个晶体学点群各自包含的对称操作。点群按七个晶系排列，每个点群的全部元素都已列出。只要点群有主轴，主轴就取在 $Oz$ 方向，这称为\textit{标准定向}（standard setting）；当然还可能有许多其他定向。图 1.4 给出 32 个晶体学点群的极射赤面投影图；投影理论可参见例如 Phillips (1963)。在同一晶系中，含有最多对称操作的点群称为该晶系的\textit{全对称点群}（holosymmetric point group）。

\textit{点群操作的乘法}

对立方群以外的所有点群，借助于极射赤面投影图容易得到两个点群操作相继作用的结果。例如乘法 $C_{2x}C_{4z}^{+} = C_{2b}$ 见图 1.5：图中黑点表示纸面上方的点，圆圈表示纸面下方的点。点 $A$ 经操作 $C_{4z}^{-}$ 变为点 $B$；$B$ 再经操作 $C_{2x}$ 变为点 $C$。总的效果是把 $A$ 直接变为 $C$ 的操作，即 $C_{2b}$，这从图 1.1 可以看出。

\setcounter{figure}{4}
\begin{figure}[H]
\centering
\includegraphics[width=0.55\textwidth]{fig1_5.png}
\caption{点 $A$ 经 $C_{4z}^{-}$ 到 $B$、再经 $C_{2x}$ 到 $C$ 的示意（原书本图无图注文字）。}
\end{figure}

对立方点群，得到群乘法表的最容易方法是使用表 1.4(a) 给出的算符的 Jones 忠实表示符号（Jones' faithful representation symbols）。算符 $R$ 的 Jones 符号就是 $R$ 作用在向量 $(x, y, z)$ 上得到的向量；这些符号正是《X 射线晶体学国际表》第一卷表 4.3 中点式立方与六方空间群的一般等价位置坐标（Henry and Lonsdale 1965）。若 $(x, y, z)$ 是相对沿坐标轴的单位矢量 $\mathbf{i}$、$\mathbf{j}$、$\mathbf{k}$ 给出的，则例如 $C_{4z}^{+}x\mathbf{i} = x\mathbf{j}$，$C_{4z}^{+}y\mathbf{j} = -y\mathbf{i}$，$C_{4z}^{+}z\mathbf{k} = z\mathbf{k}$，于是可写 $C_{4z}^{+}(x, y, z) = (-y, x, z)$，即

\begin{equation*}
C_{4z}^{+}(x\mathbf{i} + y\mathbf{j} + z\mathbf{k}) = (-y\mathbf{i} + x\mathbf{j} + z\mathbf{k})
\tag{1.4.8}
\end{equation*}

因此 $C_{4z}^{+}$ 的 Jones 符号写作 $(\bar{y}xz)$。以行向量 $\langle \mathbf{i}, \mathbf{j}, \mathbf{k}|$ 为基时代表 $C_{4z}^{+}$ 的矩阵是

\begin{equation*}
\mathbf{D}(C_{4z}^{+}) = \begin{pmatrix} 0 & -1 & 0 \\ 1 & 0 & 0 \\ 0 & 0 & 1 \end{pmatrix}
\tag{1.4.9}
\end{equation*}

从而 Jones 符号就是这个矩阵表示的缩写。式 (1.4.9) 来自

\begin{equation*}
C_{4z}^{+}\langle \mathbf{i}, \mathbf{j}, \mathbf{k}| = \langle \mathbf{i}, \mathbf{j}, \mathbf{k}| \begin{pmatrix} 0 & -1 & 0 \\ 1 & 0 & 0 \\ 0 & 0 & 1 \end{pmatrix},
\tag{1.4.10}
\end{equation*}

这说明了 $D$ 为什么是忠实表示。这些符号的乘法规则容易由矩阵乘法导出，例如

\begin{equation*}
(zxy)(\bar{y}xz) = \begin{pmatrix} 0 & 0 & 1 \\ 1 & 0 & 0 \\ 0 & 1 & 0 \end{pmatrix} \begin{pmatrix} 0 & -1 & 0 \\ 1 & 0 & 0 \\ 0 & 0 & 1 \end{pmatrix} = \begin{pmatrix} 0 & 0 & 1 \\ 0 & -1 & 0 \\ 1 & 0 & 0 \end{pmatrix} = (z\bar{y}x).
\tag{1.4.11}
\end{equation*}

其实不必写出完整矩阵：$(zxy)(\bar{y}xz) = (z'x'y')$，其中 $(x'y'z') = (\bar{y}xz)$，于是 $(z'x'y') = (z\bar{y}x)$。式 (1.4.11) 作为表示等式读作

\begin{equation*}
\mathbf{D}(C_{31}^{+})\mathbf{D}(C_{4z}^{+}) = \mathbf{D}(C_{2c})
\tag{1.4.12}
\end{equation*}

这可由表 1.4(a) 看出：表中给出了全部 48 个立方点群操作的 Jones 忠实表示符号。由于 $D$ 是忠实表示，立即可得 $C_{31}^{+}C_{4z}^{+} = C_{2c}$。表 1.4(a) 的 Jones 符号也可用于四方、正交点群的元素，以及单斜、三斜点群的元素——只要向量 $(x, y, z)$ 相对于平行于晶体坐标轴的单位矢量 $\mathbf{i}$、$\mathbf{j}$、$\mathbf{k}$ 给出。表 1.4(b) 中类似的一组符号给出六方或三方点群的操作对向量 $(x, y, z)$ 的作用，此时 $(x, y, z)$ 相对于单位矢量 $\mathbf{i}$、$\mathbf{j}$、$\mathbf{k}$ 给出，而 $\mathbf{i}$ 与 $\mathbf{j}$ 互成 120° 且都在垂直于 $\mathbf{k}$ 的平面内。

表 1.5 与表 1.6 给出点群 432（$O$）与 622（$D_{6}$）的群乘法表（Belov 1957\textsuperscript{b}，Belov and Tarkhova 1960，Hurley 1966，Koptsik 1966）。这两个点群只含纯转动；其余只含纯转动的点群都是这两个点群之一的子群。而每个含有非纯转动操作的点群，都与某个纯转动点群简单地相关联。因此，由表 1.5 与 1.6 可以很快得到 32 个点群中任何一个的群乘法表。

\begin{table}[H]
\centering
\textbf{表 1.2}\quad \textit{七个晶系的对称元素}

\medskip
\resizebox{\textwidth}{!}{%
\begin{tabular}{|c|c|c|c|c|c|c|}
\hline
三斜 & 单斜 & 正交 & 四方 & 三方 & 六方 & 立方 \\ \hline
$\bar{1}$（$C_{i}$） & $2/m$（$C_{2h}$） & $mmm$（$D_{2h}$） & $4/mmm$（$D_{4h}$） & $\bar{3}m$（$D_{3d}$） & $6/mmm$（$D_{6h}$） & $m3m$（$O_{h}$） \\ \hline
$E, I$ & $E, I$ & $E, I$ & $E, I$ & $E, I$ & $E, I$ & $E, I$ \\
 & $C_{2z}, \sigma_{z}$ & $C_{2m}, \sigma_{m}$ & $C_{4z}^{\pm}, S_{4z}^{\mp}$ & $C_{3}^{\pm}, S_{6}^{\mp}$ & $C_{6}^{\pm}, S_{3}^{\mp}$ & $C_{2m}, \sigma_{m}$ \\
 &  &  & $C_{2m}, \sigma_{m}$ & $C'_{2i}, \sigma_{di}$ & $C_{3}^{\pm}, S_{6}^{\mp}$ & $C_{3j}^{\pm}, S_{6j}^{\mp}$ \\
 &  &  & $C_{2s}, \sigma_{ds}$ &  & $C_{2}, \sigma_{h}$ & $C_{4m}^{\pm}, S_{4m}^{\mp}$ \\
 &  &  &  &  & $C'_{2i}, \sigma_{di}$ & $C_{2p}, \sigma_{dp}$ \\
 &  &  &  &  & $C''_{2i}, \sigma_{vi}$ &  \\ \hline
\end{tabular}}

\smallskip
\parbox{0.92\textwidth}{\textit{表 1.2 的注.}
(i) 三方晶系还有一种定向，以 $C''_{2i}$ 代替 $C'_{2i}$ 作为标准定向。
(ii) 下列轴的位置见图 1.1、1.2、1.3：$m = x, y, z$；$s = a, b$；$i = 1, 2, 3$；$j = 1, 2, 3, 4$；$p = a, b, c, d, e, f$。}
\end{table}

\begin{table}[H]
\centering
\textbf{表 1.3}\quad \textit{32 个点群}

\smallskip
\begin{tabular}{|c|c|c|p{0.58\textwidth}|}
\hline
编号 & 国际记号 & Schönflies 记号 & 元素 \\ \hline
\multicolumn{4}{|l|}{\textit{三斜}} \\ \hline
1 & $1$ & $C_{1}$ & $E$ \\ \hline
2 & $\bar{1}$ & $C_{i}$ & $E$, $I$ \\ \hline
\multicolumn{4}{|l|}{\textit{单斜}} \\ \hline
3 & $2$ & $C_{2}$ & $E$, $C_{2z}$ \\ \hline
4 & $m$ & $C_{s}$, $C_{1h}$ & $E$, $\sigma_{z}$ \\ \hline
5 & $2/m$ & $C_{2h}$ & $E$, $C_{2z}$, $I$, $\sigma_{z}$ \\ \hline
\multicolumn{4}{|l|}{\textit{正交}} \\ \hline
6 & $222$ & $D_{2}$ & $E$, $C_{2x}$, $C_{2y}$, $C_{2z}$ \\ \hline
7 & $mm2$ & $C_{2v}$ & $E$, $C_{2z}$, $\sigma_{x}$, $\sigma_{y}$ \\ \hline
8 & $mmm$ & $D_{2h}$ & $E$, $C_{2x}$, $C_{2y}$, $C_{2z}$, $I$, $\sigma_{x}$, $\sigma_{y}$, $\sigma_{z}$ \\ \hline
\multicolumn{4}{|l|}{\textit{四方}} \\ \hline
9 & $4$ & $C_{4}$ & $E$, $C_{4z}^{+}$, $C_{4z}^{-}$, $C_{2z}$ \\ \hline
10 & $\bar{4}$ & $S_{4}$ & $E$, $S_{4z}^{+}$, $S_{4z}^{-}$, $C_{2z}$ \\ \hline
11 & $4/m$ & $C_{4h}$ & $E$, $C_{4z}^{+}$, $C_{4z}^{-}$, $C_{2z}$, $I$, $S_{4z}^{-}$, $S_{4z}^{+}$, $\sigma_{z}$ \\ \hline
12 & $422$ & $D_{4}$ & $E$, $C_{4z}^{+}$, $C_{4z}^{-}$, $C_{2z}$, $C_{2x}$, $C_{2y}$, $C_{2a}$, $C_{2b}$ \\ \hline
13 & $4mm$ & $C_{4v}$ & $E$, $C_{4z}^{+}$, $C_{4z}^{-}$, $C_{2z}$, $\sigma_{x}$, $\sigma_{y}$, $\sigma_{da}$, $\sigma_{db}$ \\ \hline
14 & $\bar{4}2m$ & $D_{2d}$ & $E$, $S_{4z}^{+}$, $S_{4z}^{-}$, $C_{2z}$, $C_{2x}$, $C_{2y}$, $\sigma_{da}$, $\sigma_{db}$ \\ \hline
15 & $4/mmm$ & $D_{4h}$ & $E$, $C_{4z}^{+}$, $C_{4z}^{-}$, $C_{2z}$, $C_{2x}$, $C_{2y}$, $C_{2a}$, $C_{2b}$, $I$, $S_{4z}^{-}$, $S_{4z}^{+}$, $\sigma_{z}$, $\sigma_{x}$, $\sigma_{y}$, $\sigma_{da}$, $\sigma_{db}$ \\ \hline
\multicolumn{4}{|l|}{\textit{三方}} \\ \hline
16 & $3$ & $C_{3}$ & $E$, $C_{3}^{+}$, $C_{3}^{-}$ \\ \hline
17 & $\bar{3}$ & $C_{3i}$ & $E$, $C_{3}^{+}$, $C_{3}^{-}$, $I$, $S_{6}^{-}$, $S_{6}^{+}$ \\ \hline
18 & $32$ & $D_{3}$ & $E$, $C_{3}^{+}$, $C_{3}^{-}$, $C'_{21}$, $C'_{22}$, $C'_{23}$ \\ \hline
19 & $3m$ & $C_{3v}$ & $E$, $C_{3}^{+}$, $C_{3}^{-}$, $\sigma_{d1}$, $\sigma_{d2}$, $\sigma_{d3}$ \\ \hline
20 & $\bar{3}m$ & $D_{3d}$ & $E$, $C_{3}^{+}$, $C_{3}^{-}$, $C'_{21}$, $C'_{22}$, $C'_{23}$, $I$, $S_{6}^{-}$, $S_{6}^{+}$, $\sigma_{d1}$, $\sigma_{d2}$, $\sigma_{d3}$ \\ \hline
\multicolumn{4}{|l|}{\textit{六方}} \\ \hline
21 & $6$ & $C_{6}$ & $E$, $C_{6}^{+}$, $C_{6}^{-}$, $C_{3}^{+}$, $C_{3}^{-}$, $C_{2}$ \\ \hline
22 & $\bar{6}$ & $C_{3h}$ & $E$, $S_{3}^{-}$, $S_{3}^{+}$, $C_{3}^{+}$, $C_{3}^{-}$, $\sigma_{h}$ \\ \hline
23 & $6/m$ & $C_{6h}$ & $E$, $C_{6}^{+}$, $C_{6}^{-}$, $C_{3}^{+}$, $C_{3}^{-}$, $C_{2}$, $I$, $S_{3}^{-}$, $S_{3}^{+}$, $S_{6}^{-}$, $S_{6}^{+}$, $\sigma_{h}$ \\ \hline
24 & $622$ & $D_{6}$ & $E$, $C_{6}^{+}$, $C_{6}^{-}$, $C_{3}^{+}$, $C_{3}^{-}$, $C_{2}$, $C'_{21}$, $C'_{22}$, $C'_{23}$, $C''_{21}$, $C''_{22}$, $C''_{23}$ \\ \hline
25 & $6mm$ & $C_{6v}$ & $E$, $C_{6}^{+}$, $C_{6}^{-}$, $C_{3}^{+}$, $C_{3}^{-}$, $C_{2}$, $\sigma_{d1}$, $\sigma_{d2}$, $\sigma_{d3}$, $\sigma_{v1}$, $\sigma_{v2}$, $\sigma_{v3}$ \\ \hline
26 & $\bar{6}2m$ & $D_{3h}$ & $E$, $S_{3}^{-}$, $S_{3}^{+}$, $C_{3}^{+}$, $C_{3}^{-}$, $\sigma_{h}$, $C'_{21}$, $C'_{22}$, $C'_{23}$, $\sigma_{v1}$, $\sigma_{v2}$, $\sigma_{v3}$ \\ \hline
27 & $6/mmm$ & $D_{6h}$ & $E$, $C_{6}^{+}$, $C_{6}^{-}$, $C_{3}^{+}$, $C_{3}^{-}$, $C_{2}$, $C'_{21}$, $C'_{22}$, $C'_{23}$, $C''_{21}$, $C''_{22}$, $C''_{23}$, $I$, $S_{3}^{-}$, $S_{3}^{+}$, $S_{6}^{-}$, $S_{6}^{+}$, $\sigma_{h}$, $\sigma_{d1}$, $\sigma_{d2}$, $\sigma_{d3}$, $\sigma_{v1}$, $\sigma_{v2}$, $\sigma_{v3}$ \\ \hline
\multicolumn{4}{|l|}{\textit{立方}} \\ \hline
28 & $23$ & $T$ & $E$, $C_{2m}$, $C_{3j}^{+}$, $C_{3j}^{-}$ \\ \hline
29 & $m3$ & $T_{h}$ & $E$, $C_{2m}$, $C_{3j}^{+}$, $C_{3j}^{-}$, $I$, $\sigma_{m}$, $S_{6j}^{-}$, $S_{6j}^{+}$ \\ \hline
30 & $432$ & $O$ & $E$, $C_{2m}$, $C_{3j}^{+}$, $C_{3j}^{-}$, $C_{2p}$, $C_{4m}^{+}$, $C_{4m}^{-}$ \\ \hline
31 & $\bar{4}3m$ & $T_{d}$ & $E$, $C_{2m}$, $C_{3j}^{+}$, $C_{3j}^{-}$, $\sigma_{dp}$, $S_{4m}^{-}$, $S_{4m}^{+}$ \\ \hline
32 & $m3m$ & $O_{h}$ & $E$, $C_{2m}$, $C_{3j}^{+}$, $C_{3j}^{-}$, $C_{2p}$, $C_{4m}^{+}$, $C_{4m}^{-}$, $I$, $\sigma_{m}$, $S_{6j}^{-}$, $S_{6j}^{+}$, $\sigma_{dp}$, $S_{4m}^{-}$, $S_{4m}^{+}$ \\ \hline
\end{tabular}

\end{table}

\noindent\textit{表 1.3 的注.}\ （i）第 1 列的编号是任意的，采用 Koster, Dimmock, Wheeler, and Statz (1963) 的编号。（ii）对称操作的标记可由图 1.1--1.3 辨认：$j = 1, 2, 3, 4$；$m = x, y, z$；$p = a, b, c, d, e, f$。（iii）主轴已取在 $Oz$ 方向，但某些点群仍可能有其他定向，例如 $\bar{4}2m$（$D_{2d}$）也可以包含元素 $E$、$S_{4z}^{-}$、$C_{2z}$、$S_{4z}^{+}$、$\sigma_{x}$、$\sigma_{y}$、$C_{2a}$ 与 $C_{2b}$；在考虑空间群（见第 3 章）时这类不同的取法有时很重要。

\setcounter{figure}{3}
\begin{figure}[H]
\centering
\includegraphics[width=0.88\textwidth]{fig1_4a_scan.png}
\end{figure}
\clearpage
\begin{figure}[H]
\centering
\includegraphics[width=0.88\textwidth]{fig1_4b_scan.png}
\caption{32 个点群的一般等价方向极点的极射赤面投影图及各点群的对称元素（所有图中 $z$ 轴均垂直于纸面；$m3$ 与 $m3m$ 只画出一半的点）（Henry and Lonsdale 1965）。列：三斜、单斜（第一种定向）、四方、三方、六方、立方，以及单斜（第二种定向）与正交；行：$\chi$；$\bar{\chi}$（偶）；$\chi$（偶）加对称中心与 $\chi$（奇）；$\chi 2$；$\chi m$；$\bar{\chi}2$（偶）或 $\bar{\chi}m$（偶）；$\chi 2$ 或 $\chi m$ 加对称中心与 $\chi m$（奇）。}
\end{figure}
\setcounter{figure}{5}
